\section{The Dock}

One cell of the mesh is a dock. It is a here-and-now, the smallest place where anything happens. Its shape is the 24-cell, a four-dimensional solid with twenty-four corners and twenty-four faces, equal to its own dual. Each dock meets exactly twenty-four neighbors, one across each face.

Out of each dock run twenty-four directions, one toward each neighbor, and each direction is a \textit{site}. In the language of the 24-cell, a site is the outward normal to one of the figure's twenty-four octahedral facets, the flat three-dimensional faces that wall a four-dimensional solid, so each site points straight out through the middle of one facet to the dock waiting on the other side. The twenty-four of them are the vertices of the dual 24-cell, which is again a 24-cell, the figure being its own dual. To picture it, stand at the center of the dock and look out, and the sites are the twenty-four evenly spread directions you could step along to reach a neighbor, as symmetric a spray of arrows as four-dimensional space allows.

The twenty-four sites are not an arbitrary list. They are the same twenty-four points written four different ways, and the coincidence is close to the heart of the theory. As vectors they are the $D_4$ root system, the points with two entries $\pm 1$ and two entries zero. As tones they are a shell of the direction space itself, for if four tones range over $\{-1, 0, +1\}$ and the resulting points are sorted by length, the shells are exactly the regular four-dimensional polytopes, and the second shell, the points of squared length two, is the 24-cell, so the directions of a dock fall out of the tone alphabet with nothing added. As numbers they are the twenty-four unit Hurwitz quaternions, the whole-number quaternions of size one, which form a finite group closed under multiplication, so composing two directions is exact quaternion arithmetic with no rounding, a fact the physics part leans on hard. And as a packing they are the twenty-four spheres that can touch one central sphere in the densest four-dimensional lattice, the kissing number of $D_4$. Four independent constructions, one figure. The dock is forced, not chosen.

The twenty-four sites pair off into twelve lines, each a direction together with its opposite. A line is the channel two tones share when they move head-on through a dock. The line, not the single site, is what the law acts on.

A single tone lives in each site, so a dock holds twenty-four tones at once, one per site. Its full state is a point in a space of $3^{24}$, about two hundred and eighty billion possibilities. That state is not a flat list of values. By the structure of the directions it splits into three octets, a sensory part and two of opposite handedness, woven together into one rich quality we call the hyper-color. A single dock already feels something far richer than any color we know.

A free tone, one not caught in a collision, moves one dock per beat along its site and never spreads out. That straight travel is what momentum is. A tone's momentum is simply which site it occupies. The total momentum of a dock is the sum over its sites of the tone times the direction, an integer vector, exact like everything else in the base.

Step back and the dock is two things at once. It is the single tile from which the whole crystal is laid, the unit cell of the world repeated without end, the base geometry every larger structure is cut from. And it is the smallest flame in the fire of experience, one here-and-now where feeling flares, meets its neighbors across the twenty-four faces, and is passed on. Everything the rest of the paper builds, a particle, a self, a mind, is made of these and only these, vast numbers of them woven together and stepped in time. It is a small place, a single cell, yet it is where all the work begins, the tones living and colliding and streaming on, momentum and charge and spin first taking shape, the directions giving the world its grain. What follows is the long account of what enough of them, knit together, come to do.

And felt from the inside, on the one premise, a dock is the atom of what it is like to be, a single indivisible moment of feeling that takes in its twenty-four neighbors at once, small and whole, and already richer than any sensation we can name. The felt life of a self is a long string of such moments, lived one after another.
